\chapter{Introduction}
\label{chapter:introduction}

The living cell is in constant motion.
Proteins diffuse, assemble, and disassemble; organelles reshape and migrate;
signalling cascades propagate on timescales of seconds to minutes.
Capturing these dynamic processes demand simultaneous high spatial and temporal
resolution, a combination that conventional optical microscopy cannot simultaneously provide.
The diffraction limit of light restricts spatial resolution to roughly
200--300\,nm~\cite{born1999principles}, far larger than the nanometre-scale
structures that govern cellular function. These structures are entirely hidden within the
diffraction blur.
% Fluorescent microscopy utilizes fluorescent molecules (fluorophores) that are attached
% to imaged structure. Upon excitation with a laser blinking fluorophores are captured
% in a sequence of images and reconstructed.
% Different fluorophores can emit light at different intervals making them more suitable for 
% certain reconstruction methods. 


Fluorescence microscopy addresses this challenge by labelling structures of 
interest with fluorescent molecules (fluorophores). Upon excitation with a laser, 
fluorophores emit light that can be detected and imaged. 
Several super-resolution microscopy techniques have been developed to surpass 
the diffraction limit and recover sub-diffraction information. 
However, these methods typically involve significant trade-offs between 
spatial resolution, acquisition speed, phototoxicity, hardware complexity, 
and computational cost.

% Among these trade-offs, temporal resolution is often the first to be
% sacrificed, yet it is precisely what live-cell imaging cannot do without.
% Many of the most important cellular events are transient and unpredictable:
% vesicle fusion, protein recruitment to a damage site, or the brief
% reorganisation of the cytoskeleton can occur within seconds and never repeat in
% the same way. Methods that reconstruct an image offline from a long acquisition
% force the observer to commit to a fixed time window before knowing what will
% happen, so rare or fast events are easily missed or averaged away. Real-time
% super-resolution, by contrast, lets the experimenter watch dynamics as they
% unfold and intervene during live experiments---for example, following a moving
% organelle or triggering a perturbation at the right moment. This closes the
% loop between observation and control, turning microscopy from a record of what
% already happened into a tool for interacting with living systems.

Structured Illumination Microscopy (SIM)~\cite{sim2000,gustafsson2000sim} 
and Stimulated Emission Depletion (STED)~\cite{sted1994} 
achieve super-resolution through specialised optical hardware and 
illumination strategies.
Single-Molecule Localisation Microscopy (SMLM) methods \cite{betzig_2006_imaging,rust_2006_subdiffractionlimit}, 
achieve even higher resolution by localising individual fluorophores. 
However, SMLM requires sparse stochastic blinking and typically relies on 
tens of thousands of frames to reconstruct a single super-resolved image. 
Acquisition times can therefore range from minutes to hours, 
severely limiting temporal resolution and precluding real-time observation of 
dynamic cellular processes.

Fluctuation-based methods (Super-Resolution Optical Fluctuation Imaging
(SOFI)~\cite{sofi2009} and Super-Resolution Radial Fluctuations
(SRRF)~\cite{srrf2016}) offer a more favourable balance between spatial and temporal resolution.
These methods operate with dense fluorophore labelling and can achieve up to n-fold spatial improvement using higher order statistics.
 
Despite these advantages, SOFI still requires hundreds of frames per
reconstruction.
Each SR image is computed offline from the entire acquired sequence, with
computational latency in the range of seconds even with GPU
acceleration~\cite{resurf2024}.
This is incompatible with real-time live-cell imaging.

Deep learning offers a promising route to reduce both the number of frames
required and the inference latency.
Rather than computing cumulants analytically, a neural network can learn to map
a short sequence of low-resolution (LR) fluorescence frames directly to a
second-order SR image, effectively approximating the statistical operation from
far fewer observations.

Tekpinar et al.~\cite{resurf2024} introduced
SOFI-MISRGRU, adapting the Multi-Image Super-Resolution GRU
(MISRGRU) architecture~\cite{MISRGRU-original} (originally developed for
satellite imagery) to the fluorescence microscopy domain.
The model demonstrated a
400-fold latency reduction relative to conventional SOFI establishing the
practical feasibility of real-time deep-learning SR microscopy.

In this thesis we explore whether transformer-based fusion module (TR-MISR) \cite{TR-MISR-original}
can match or surpass the SR quality of the SOFI-MISRGRU \cite{resurf2024} baseline on simulated
fluorescence microscopy data across varying SNR conditions. On PROBA-V dataset TR-MISR achieved better results and 
demonstrated robustness to noisy images a property the authors attribute to the attention mechanism.

Secondly, we also extend SOFI-MISRGRU into a continuous-output streaming architecture 
that eliminates the batch acquisition bottleneck entirely. With this design the model
can combine more input images which enchances the reconstruction quality while 
enabling uninterrupted super-resolution video generation from dynamic samples in real-time.


% The main contributions of this thesis are:

% \begin{enumerate}
% \item Adaptation and evaluation of the TR-MISR transformer-based fusion
%       architecture for fluorescence microscopy super-resolution.
% \item Systematic comparison of TR-MISR and SOFI-MISRGRU across multiple SNR
%       conditions on simulated fluorescence microscopy data.
% \item Development of a streaming SOFI-MISRGRU architecture that produces one
%       super-resolved output per input frame, enabling continuous reconstruction
%       without fixed-window batch constraints,.
% \end{enumerate}

\section{Thesis Outline}

Chapter~\ref{ch:background} provides background on fluorescence microscopy, SOFI
theory, and the deep-learning building blocks underlying both architectures.
Chapter~\ref{ch:method} describes the adaptations made to TR-MISR and MISRGRU
for this application, including the encoder, Transformer fusion module, decoder
variants, and the motion augmentation training strategy.
Chapter~\ref{ch:datasets} details the synthetic simulation pipeline and
training and test set construction.
Chapter~\ref{ch:results} presents the experimental results, addressing each
research question in turn.
Chapter~\ref{ch:discussion} interprets the findings in the context of real-time
live-cell imaging and identifies limitations.

